% Standalone problem document, generated from the adjacent README.md.
% Compile with XeLaTeX. Mathematical content is maintained in Markdown.
\documentclass[11pt,a4paper]{article}
\usepackage[fontsize=10.5pt]{scrextend}
\usepackage[margin=22mm,headheight=15pt,headsep=9mm,footskip=11mm]{geometry}
\usepackage{fontspec}
\setmainfont{texgyrepagella}[Extension=.otf,UprightFont=*-regular,BoldFont=*-bold,ItalicFont=*-italic,BoldItalicFont=*-bolditalic]
\setsansfont{texgyreheros}[Extension=.otf,UprightFont=*-regular,BoldFont=*-bold,ItalicFont=*-italic,BoldItalicFont=*-bolditalic]
\setmonofont{texgyrecursor}[Extension=.otf,UprightFont=*-regular,BoldFont=*-bold,ItalicFont=*-italic,BoldItalicFont=*-bolditalic,Scale=MatchLowercase]
\usepackage{amsmath,amssymb}
\usepackage{unicode-math}
\setmathfont{texgyrepagella-math.otf}
\usepackage{xcolor,microtype,enumitem,needspace,fancyhdr}
\usepackage{bookmark}
\definecolor{ink}{HTML}{17344B}
\definecolor{accent}{HTML}{197D80}
\definecolor{muted}{HTML}{596773}
\hypersetup{colorlinks=true,linkcolor=accent,urlcolor=accent,pdftitle={RA-07 - Convexity
of the expected error of volume-sampled column
subsets},pdfauthor={Open Problems in Numerical Linear Algebra}}
\urlstyle{same}
\setlength{\parindent}{0pt}
\setlength{\parskip}{5pt plus 1pt minus 1pt}
\linespread{1.04}
\setlength{\emergencystretch}{3em}
\widowpenalty=10000
\clubpenalty=10000
\displaywidowpenalty=10000
\allowdisplaybreaks[1]
\setlist{nosep,leftmargin=1.5em}
\providecommand{\tightlist}{\setlength{\itemsep}{0pt}\setlength{\parskip}{0pt}}
\providecommand{\pandocbounded}[1]{#1}
\pagestyle{fancy}
\fancyhf{}
\fancyhead[L]{\sffamily\footnotesize\color{muted}OPEN PROBLEMS IN NUMERICAL LINEAR ALGEBRA}
\fancyhead[R]{\sffamily\bfseries\footnotesize\color{accent}RA-07}
\fancyfoot[L]{\parbox[b]{0.9\textwidth}{\sffamily\fontsize{8}{10}\selectfont\color{muted}Editorial ratings; a bounded search cannot certify that a problem remains unsolved.\\Literature check: 2026-09-12}}
\fancyfoot[R]{\sffamily\footnotesize\color{muted}\thepage}
\renewcommand{\headrulewidth}{0.3pt}
\renewcommand{\headrule}{\hbox to\headwidth{\color{accent}\leaders\hrule height \headrulewidth\hfill}}
\makeatletter
\renewcommand{\section}{\Needspace{5\baselineskip}\@startsection{section}{1}{0pt}{12pt plus 2pt minus 2pt}{4pt}{\sffamily\bfseries\large\color{ink}}}
\renewcommand{\subsection}{\Needspace{4\baselineskip}\@startsection{subsection}{2}{0pt}{9pt plus 1pt minus 1pt}{3pt}{\sffamily\bfseries\normalsize\color{ink}}}
\makeatother
\setcounter{secnumdepth}{0}
\begin{document}
{\sffamily\small\color{muted}Randomized and low-rank approximation\par}
\vspace{3pt}
{\raggedright\hyphenpenalty=10000\exhyphenpenalty=10000\sffamily\bfseries\fontsize{21}{24}\selectfont\color{ink}Convexity
of the expected error of volume-sampled column subsets\par}
\vspace{7pt}
{\sffamily\small\textbf{Difficulty:} challenging\quad\textbf{Importance:} interesting
to the community\par}
{\sffamily\small\bfseries\color{accent}Status: Lean verified\par}
\vspace{4pt}
{\color{accent}\hrule height 0.8pt}
\vspace{6pt}
\textbf{Rating rationale:} Challenging because the required
second-difference inequality goes beyond standard log-concavity;
community impact is understanding average column-selection and Nyström
error.

\subsection{Lean proof and verification evidence ---
2026-09-12}\label{lean-proof-and-verification-evidence-2026-09-12}

\textbf{The complete original discrete-convexity assertion is Lean
verified.} The
\href{https://github.com/sgstepaniants/OpenProblemsInNLA/tree/bf144a8ea84992d64f79f4425b18352843376286/randomized-and-low-rank-approximation/RA-07/lean}{proof
at revision bf144a8} covers every \(n\geq3\), strictly positive real
tuple and \(2\leq j\leq n-1\), including repeated eigenvalues and both
endpoint indices, without ordering or factorization assumptions.

\textbf{Mathematical proof:} Matthew J. Colbrook, Department of Applied
Mathematics and Theoretical Physics, University of Cambridge.
\textbf{Lean formalization:} George Stepaniants, Department of Computing
and Mathematical Sciences, California Institute of Technology, Pasadena,
California, USA, with AI-agent assistance.

The actual generating-polynomial derivative of order \(j-1\) has
\(n-j+1\geq2\) positive reciprocal-root factors \(\mu_a\), with
multiplicity. The proof derives

\[f(j-1)-2f(j)+f(j+1)
=\frac{2\sum_{a< b}\mu_a\mu_b(\mu_a-\mu_b)^2}
{s_1(s_1^2-s_2)}\geq0,
\qquad s_r=\sum_a\mu_a^r,\]

including strict positivity of the denominator. The six
\href{https://github.com/sgstepaniants/OpenProblemsInNLA/blob/bf144a8ea84992d64f79f4425b18352843376286/randomized-and-low-rank-approximation/RA-07/lean/Solution.lean}{checked
exports}, each with prefix \textit{NLA.RA07.}, are:

\begin{itemize}
\tightlist
\item
  \textit{elementary\_values}: actual subset-sum conventions and
  positivity.
\item
  \textit{generating\_derivative\_values}: actual coefficients and
  factorial-scaled derivatives.
\item
  \textit{positive\_derivative\_factorization}: exact degree, real
  splitting and every root multiplicity.
\item
  \textit{power\_sum\_certificate}: positive denominator and exact
  nonnegative pair identity.
\item
  \textit{second\_difference\_certificate}: that identity for every
  original ratio and index.
\item
  \textit{errorSequence\_convex}: the complete affirmative canonical
  theorem.
\end{itemize}

Separate pairs of independent agents approved the
\href{https://github.com/ajt60gaibb/OpenProblemsInNLA/blob/main/randomized-and-low-rank-approximation/RA-07/lean/reviews}{statements
and proof}.
\href{https://github.com/sgstepaniants/OpenProblemsInNLA/actions/runs/34715563781}{Linux
run 34715563781} matched all six exports with the sandboxed Comparator
and replayed the solution in Lean's default kernel. The
\href{https://github.com/ajt60gaibb/OpenProblemsInNLA/blob/main/randomized-and-low-rank-approximation/RA-07/lean/verification/linux-2026-09-12}{operational
audit and original artifacts} bind 123 inputs and the
isolation/rejection controls. Twelve transitive axiom checks use only
\textit{propext}, \textit{Classical.choice} and \textit{Quot.sound}. No
external human peer review is claimed.

The pins are \textbf{Lean 4.33.1},
\href{https://github.com/alerad/leancert/tree/621a43d7cf21f87872392a01e874f2f1dbddc926}{LeanCert
621a43d} and
\href{https://github.com/leanprover-community/mathlib4/tree/0df444a360eaa60ab8c11dca51a86af692955474}{Mathlib
0df444a}. LeanCert audits kernel trust; there is no numerical interval
certificate or approximate root computation. See the
\href{https://github.com/ajt60gaibb/OpenProblemsInNLA/blob/main/randomized-and-low-rank-approximation/RA-07/lean/formalization.yaml}{manifest},
\href{https://github.com/ajt60gaibb/OpenProblemsInNLA/blob/main/randomized-and-low-rank-approximation/RA-07/lean/lake-manifest.json}{dependency
pins} and
\href{https://github.com/ajt60gaibb/OpenProblemsInNLA/blob/main/randomized-and-low-rank-approximation/RA-07/lean/NUMERICAL_TARGETS.md}{exact
targets}. From the verified revision on a
\href{https://github.com/ajt60gaibb/OpenProblemsInNLA/blob/main/tools/lean/HARNESS.md}{non-root
Linux host}:

\begin{verbatim}
tools/lean/bootstrap.sh /absolute/path/to/nla-lean-tools
tools/lean/selftest.sh /absolute/path/to/nla-lean-tools
tools/lean/verify.sh \
  randomized-and-low-rank-approximation/RA-07/lean \
  /absolute/path/to/nla-lean-tools
\end{verbatim}

The source's additional monotonicity, index-one, sampling-expectation
and application claims are outside these six exports. The complete
original scalar convexity target and its historical informal proof
remain below.

\subsection{Resolution --- 2026-09-11}\label{resolution-2026-09-11}

\textbf{Author:} Matthew J. Colbrook, Department of Applied Mathematics
and Theoretical Physics, University of Cambridge. \textbf{Independent
Codex-agent review: PASS for the exact target.}

The sequence \((j+1)e_{j+1}/e_j\) is decreasing and discretely convex
for every positive spectrum, including both endpoints. The exact
second-difference certificate proves the full canonical conjecture.

The complete target is resolved. Its former difficulty rating is
historical; the original mathematical statement and source evidence are
retained below.

\textbf{Primary reference:}
\href{https://github.com/ajt60gaibb/OpenProblemsInNLA/blob/main/references/colbrook-transfer-2026-09-11/manuscripts/01_volume_sampling_convexity.pdf}{complete
authored PDF},
\href{https://github.com/ajt60gaibb/OpenProblemsInNLA/blob/main/references/colbrook-transfer-2026-09-11/manuscripts/01_volume_sampling_convexity.tex}{standalone
TeX}, \textbf{Theorem 1.1}.
\href{https://github.com/ajt60gaibb/OpenProblemsInNLA/blob/main/references/colbrook-transfer-2026-09-11/verification/reviews/RA-07-review.md}{Independent
proof review} ·
\href{https://github.com/ajt60gaibb/OpenProblemsInNLA/blob/main/references/colbrook-transfer-2026-09-11/README.md}{Authorship
and submission record}. The 2026-09-11 verification was independent
agent review, without external human peer review or formal
certification. The later Lean verification above covers the complete
canonical scalar convexity target.

For \(n\geq3\) and positive real numbers \(\lambda_1,\ldots,\lambda_n\),
let

\[e_j(\lambda)=\sum_{\substack{S\subseteq\{1,\ldots,n\}\\|S|=j}}
\prod_{i\in S}\lambda_i,
\qquad e_0=1,\quad e_{n+1}=0.\]

Define \(f(j)=(j+1)e_{j+1}(\lambda)/e_j(\lambda)\) for
\(1\leq j\leq n\). Is \(f\) discretely convex, that is,

\[f(j-1)-2f(j)+f(j+1)\geq0
\qquad(2\leq j\leq n-1),\]

for every such \(n\) and positive tuple \(\lambda\)?

For a matrix \(A\) with positive eigenvalues \(\lambda_i\) of \(A^TA\),
sample a set \(S\) of \(j\) columns with probability proportional to
\(\det(A_S^TA_S)\). Its expected squared Frobenius projection error is

\[\mathbb E\|(I-P_{A_S})A\|_F^2=f(j),\]

where \(P_{A_S}\) is the orthogonal projector onto the selected column
span. Thus the conjecture asserts a diminishing marginal improvement in
the average reconstruction error of these fixed-size determinantal
samples. It is distinct from optimizing a worst-case approximation
factor over all matrices with a prescribed spectrum.

\subsection{References}\label{references}

\begin{enumerate}
\def\labelenumi{\arabic{enumi}.}
\tightlist
\item
  M. Dereziński, R. Khanna, and M. W. Mahoney, \emph{Improved guarantees
  and a multiple-descent curve for Column Subset Selection and the
  Nyström method}, NeurIPS 2020, arXiv:2002.09073. Section 5, Conjecture
  1 is exactly the discrete-convexity statement.
  \href{https://arxiv.org/abs/2002.09073}{Paper}.
\item
  A. Deshpande, L. Rademacher, S. S. Vempala, and G. Wang, \emph{Matrix
  Approximation and Projective Clustering via Volume Sampling}, Theory
  of Computing 2 (2006), Article 12, pp.~225--247. See the
  volume-sampling matrix approximation results.
  \href{https://theoryofcomputing.org/articles/v002a012/}{Paper}.
\end{enumerate}

\subsection{Status check --- 2026-09-08}\label{status-check-2026-09-08}

Searched ``Derezinski Khanna Mahoney Conjecture 1 convexity'',
``elementary symmetric ratio \(k+1\) convexity DPP'', and the exact 2020
paper title with 2025/2026. Checked the current arXiv record. No proof
or counterexample was found. Standard Newton inequalities imply other
ratio inequalities but are not cited as a resolution of this
second-difference inequality.

\subsection{Audit --- 2026-09-10}\label{audit-2026-09-10}

Rechecked
\href{https://arxiv.org/html/2002.09073}{Dereziński--Khanna--Mahoney,
§5, Conjecture 1}; the current record remains v3 of December 2020.
Volume-sampling convexity and elementary-symmetric-ratio searches found
no resolution. This remains historical explicit conjecture evidence,
supplemented by a bounded follow-up search.
\end{document}
